\documentclass[12pt]{article}
\usepackage[utf8]{inputenc}
\usepackage[T1]{fontenc}
\usepackage{hyperref}
\title{Itneg, Masadiit Language Resources for K-12 Curriculum Development}
\author{Compiled by Claude Sonnet 5 (claude-sonnet-5) for Andr\'as Kornai}
\date{}
\begin{document}
\maketitle

\section{Introduction}
Masadiit Itneg (ISO 639-3: tis) is a Cordilleran language of northern Luzon, spoken in Abra (Sallapadan, Bucloc, Boliney) and along the western edge of Kalinga. Two dialects are named: Masadiit Boliney and Masadiit Sallapadan. Ethnologue-derived figures give about 7,500 speakers, while Joshua Project gives 11,000 for the people group. Little dedicated documentation was found.

\section{Primary Educational Resources}
No Masadiit primers, textbooks or school materials were found in search results.

\section{Linguistic and Reference Materials}
\begin{itemize}
\item Wikipedia, Itneg languages (classification, dialects, locations): \url{https://en.wikipedia.org/wiki/Itneg_languages}
\item Descriptions of Philippine Languages by Drew Dolot (author page; search did not confirm any Masadiit content): \url{https://independent.academia.edu/DrewDolot}
\item A 2025 University of the Philippines linguistics thesis archive mentions work on Itneg Inlaud, a different Itneg variety, not Masadiit: \url{https://linguistics.upd.edu.ph/wp-content/uploads/2026/06/20260629_TheArchive-6-2_Regular-2025_ThesisDissertations.pdf}
\end{itemize}

\section{Religious and Community Resources}
Joshua Project people-group profile for Masadiit Itneg: \url{https://joshuaproject.net/people_groups/13562}

\section{Digital and Audio Resources}
None found in search results.

\section{Curriculum Development Context}
The language is not known to be taught in schools. No grammar, dictionary or orthography specific to Masadiit was found, so material would have to be built from new documentation, ideally with related Itneg and Kalinga varieties and the Philippine mother-tongue-based multilingual education framework in view.
\end{document}